\begin{table}[htbp]
\centering
\caption{La tríada de cuentas de transferencias nacionales: los tres vértices}
\label{tab:C13_1}
\small
\begin{tabular}{lSSSp{3.2cm}p{4.2cm}}
\toprule
{Vértice} & {2025 (mdp)} & {\% PIB} & {Var. real \%} & {Dirección demográfica} & {¿Horizonte en el paquete?} \\
\midrule
{Pensiones contributivas (clasificacion economica)} & \num{1637665.1} & \num{4.53} & \num{4.7} & {la transicion lo empuja al alza} & {SI: el CGPE proyecta pensiones y jubilaciones a 2030 (4.5 a 4.8 \% del PIB)} \\
{Pensiones no contributivas (ramo 20)} & \num{527389.1} & \num{1.46} & \num{2.6} & {la transicion lo empuja al alza} & {NO: el paquete no proyecta el padron} \\
{Salud (funcion 2.3, GF + entidades, bruta)} & \num{888903.7} & \num{2.46} & \num{-12.2} & {al alza por envejecimiento} & {NO: ninguna proyeccion por funcion} \\
{Educacion (funcion 2.5, GF, bruta)} & \num{1084590.9} & \num{3.0} & \num{0.7} & {LA TRANSICION LO EMPUJA A LA BAJA} & {NO: el paquete no declara matricula, cobertura, planteles ni docentes} \\
\bottomrule
\end{tabular}
\par\vspace{2pt}\footnotesize Fuente: elaboración propia del ITED con los analíticos del PEF aprobado 2024 y 2025 y el CGPE 2025.
\end{table}
